\documentclass{beamer}
\usepackage[utf8]{inputenc}

\usetheme{Madrid}
\usecolortheme{default}
\usepackage{amsmath,amssymb,amsfonts,amsthm}
\usepackage{txfonts}
\usepackage{tkz-euclide}
\usepackage{listings}
\usepackage{adjustbox}
\usepackage{array}
\usepackage{tabularx}
\usepackage{gvv}
\usepackage{lmodern}
\usepackage{circuitikz}
\usepackage{tikz}
\usepackage{graphicx}
\usepackage{mathtools}
\setbeamertemplate{page number in head/foot}[totalframenumber]

\usepackage{tcolorbox}
\tcbuselibrary{minted,breakable,xparse,skins}



\definecolor{bg}{gray}{0.95}
\DeclareTCBListing{mintedbox}{O{}m!O{}}{%
  breakable=true,
  listing engine=minted,
  listing only,
  minted language=#2,
  minted style=default,
  minted options={%
    linenos,
    gobble=0,
    breaklines=true,
    breakafter=,,
    fontsize=\small,
    numbersep=8pt,
    #1},
  boxsep=0pt,
  left skip=0pt,
  right skip=0pt,
  left=25pt,
  right=0pt,
  top=3pt,
  bottom=3pt,
  arc=5pt,
  leftrule=0pt,
  rightrule=0pt,
  bottomrule=2pt,
  toprule=2pt,
  colback=bg,
  colframe=orange!70,
  enhanced,
  overlay={%
    \begin{tcbclipinterior}
    \fill[orange!20!white] (frame.south west) rectangle ([xshift=20pt]frame.north west);
    \end{tcbclipinterior}},
  #3,
}
\lstset{
    language=C,
    basicstyle=\ttfamily\small,
    keywordstyle=\color{blue},
    stringstyle=\color{orange},
    commentstyle=\color{green!60!black},
    numbers=left,
    numberstyle=\tiny\color{gray},
    breaklines=true,
    showstringspaces=false,
}

\title 
{5.2.64}
\date{October 12, 2025}


\author 
{Vivek K Kumar - EE25BTECH11062}



\begin{document}


\frame{\titlepage}
\begin{frame}{Question}
Solve for the following system of linear equations:
\begin{align*}
\vec{X} + \vec{Y} = \myvec{7 & 0 \\ 2 & 5}\\
\vec{X} - \vec{Y} = \myvec{3 & 0 \\ 0 & 4}\\
\end{align*}
\end{frame}

\begin{frame}{Solution}
The given system can be represented as:
\begin{align}
    \myvec{\vec{I} & \vec{I} \\ \vec{I} & -\vec{I}}\myvec{\vec{X} \\ \vec{Y}} = \myvec{7 & 0 \\ 2 & 5 \\ 3& 0 \\ 0 & 4} 
\end{align}
Solving the equation,
\begin{align}
    \myvec{\vec{X} \\ \vec{Y}} = \myvec{1 & 0 & 1 & 0 \\ 0 & 1 & 0 & 1 \\ 1 & 0 & -1 & 0 \\ 0 & 1 & 0 & -1} ^{-1} \myvec{7 & 0 \\ 2 & 5 \\ 3& 0 \\ 0 & 4} \\
    \myvec{\vec{X} \\ \vec{Y}} = \frac{1}{2}\myvec{1 & 0 & 1 & 0 \\ 0 & 1 & 0 & 1 \\ 1 & 0 & -1 & 0 \\ 0 & 1 & 0 & -1} \myvec{7 & 0 \\ 2 & 5 \\ 3& 0 \\ 0 & 4}
\end{align}
\end{frame}
\begin{frame}{Solution}
\begin{align}
     \myvec{\vec{X} \\ \vec{Y}} = \myvec{5 & 0 \\ 1 & 9/2 \\ 2 & 0 \\ 1 & 1/2}
\end{align}
Hence, the matrices are given by
\begin{align}
    \vec{X} = \myvec{5 & 0 \\ 1 & 9/2}\\
    \vec{Y} = \myvec{2 & 0 \\ 1 & 1/2}
\end{align}
\end{frame}

\begin{frame}{Results}
\begin{table}[H]    
  \centering
  \begin{tabular}[12pt]{ |c| c|}
    \hline
    \textbf{Matrix} & \textbf{Value}\\ 
    \hline
    $\vec{X}$ & $\myvec{5 & 0 \\ 1 & 9/2}$\\
    \hline
    $\vec{Y}$ & $\myvec{2 & 0 \\ 1 & 1/2}$\\
    \hline
\end{tabular}
  \caption{Results}
  \label{tab:5.2.64}
\end{table}

\end{frame}
\end{document}